\section{Ход работы}

\subsection{Юстировка измерительной установки}

Прежде всего мы юстировали установку. Для этого оправу с излучателем аккуратно приподняли, чтобы луч едва касался прозрачной крышки столика. На её шероховатой поверхности чётко отобразился след луча, что позволило визуально контролировать его траекторию.

Далее, совмещая два действия — плавный поворот столика и регулировку юстировочным винтом — мы добились, чтобы световое пятно проходило строго над центром отверстия в столике и совпадало с отметкой 180°. Как только необходимое положение было достигнуто, оба фиксирующих винта на рейтере надёжно закрепили.

\subsection{Измерение преломляющего угла призмы}  
После юстировки установки провели серию измерений углов отражения луча от граней призмы. Столик с призмой вращали, фиксируя координаты отражённого луча по лимбу от длинного катета и гипотенузы в диапазоне \(10^\circ\)–\(140^\circ\) с шагом \(10^\circ\). Результаты измерений отсчетов по риске (стрелочке) занесены в таблицу \ref{tab:prism-angle}.

% Добавить в преамбулу: \usepackage{longtable}
\begin{longtable}[c]{|c|c|c|}
\caption{Результаты измерений углов отражения}
\label{tab:prism-angle}\\
\hline
\textbf{Координата луча, \(^\circ\)} & \textbf{Длинный катет, \(^\circ\)} & \textbf{Гипотенуза, \(^\circ\)} \\ \hline 
\endfirsthead
%
\endhead
%
10 & 64.5 & 282.0 \\ \hline
20 & 69.4 & 287.8 \\ \hline
30 & 74.5 & 292.8 \\ \hline
40 & 79.2 & 298.0 \\ \hline
50 & 84.4 & 303.0 \\ \hline
60 & 89.4 & 307.5 \\ \hline
70 & 94.5 & 312.0 \\ \hline
80 & 99.5 & 318.0 \\ \hline
90 & 104.0 & 323.0 \\ \hline
100 & 109.0 & 328.0 \\ \hline
110 & 114.2 & 333.0 \\ \hline
120 & 120.0 & 338.0 \\ \hline
130 & 124.0 & 343.0 \\ \hline
140 & 129.0 & 347.5 \\ \hline
\end{longtable}

Программная обработка данных позволила определить:  
\begin{itemize}  
    \item Преломляющий угол призмы: \(A = \left(38.43 \pm 0.19\right)^\circ\), \((\varepsilon = 0.5\%)\).  
    \item Среднеквадратическая погрешность: \(\sigma = 0.5^\circ\).  
\end{itemize}  

\begin{figure}[h]  
\centering  
\includegraphics[width=0.7\textwidth]{prism_angle_plot}  
\caption{Преломляющий угол призмы}  
\label{fig:prism-angle-plot}  
\end{figure}  

\subsection{Определение обыкновенной и необыкновенной волн}  
Для разделения лучей на обыкновенную и необыкновенную волны использовали поляроид, регулируя его ориентацию. Ключевые этапы эксперимента:  

\begin{longtable}[c]{|c|c|c|}
\caption{Зависимость угла отклонения луча от ориентации поляроида}
\label{tab:polarization}\\
\hline
\textbf{Ориентация поляроида} & \textbf{Угол поворота, \(^\circ\)} & \textbf{Тип волны} \\ \hline 
\endfirsthead
%
\endhead
%
Вертикальная & 300 & Обыкновенная \\ \hline
Горизонтальная & 200 & Необыкновенная \\ \hline
\end{longtable}

\begin{itemize}
    \item При повороте поляроида на \(300^\circ\) (вертикальная ориентация разрешённого направления) регистрировался \textbf{обыкновенный луч} с углом отклонения \(207.2^\circ\).
    \item При повороте на \(200^\circ\) (горизонтальная ориентация) наблюдался \textbf{необыкновенный луч}, отклоняющийся на \(200.0^\circ\).
\end{itemize}

Это подтверждает, что обыкновенная волна сильнее отклоняется от первоначального направления, что характерно для анизотропных сред. Разница в углах (\(207.2^\circ - 200.0^\circ = 7.2^\circ\)) обусловлена различием показателей преломления для разных поляризаций.

\subsection{Измерение главных показателей преломления}  
Для определения показателей преломления обыкновенной (\(n_o\)) и необыкновенной (\(n_e\)) волн провели серию измерений углов преломления при изменении координаты отражённого луча от входной грани призмы. Диапазон углов составил \(20^\circ\)–\(140^\circ\) с шагом \(10^\circ\). Результаты занесены в таблицу \ref{tab:refraction-data}.

% Добавить в преамбулу: \usepackage{longtable}
\begin{longtable}[c]{|c|c|c|c|}
\caption{Результаты измерений углов отражения и преломления}
\label{tab:refraction-data}\\
\hline
\textbf{Отраж.\newline (лимб), \(^\circ\)} & \textbf{Отраж.\newline (риска), \(^\circ\)} & \textbf{Обыкн.\newline, \(^\circ\)} & \textbf{Необыкн.\newline, \(^\circ\)} \\ \hline 
\endfirsthead
%
\multicolumn{4}{c}{{\tablename} \thetable\ — продолжение} \\
\hline
\textbf{Отраж.\newline (лимб), \(^\circ\)} & \textbf{Отраж.\newline (риска), \(^\circ\)} & \textbf{Обыкн.\newline, \(^\circ\)} & \textbf{Необыкн.\newline, \(^\circ\)} \\ \hline 
\endhead
%
20 & 69.3 & 211.5 & 201.8 \\ \hline
30 & 74.5 & 209.2 & 200.8 \\ \hline
40 & 79.3 & 208.0 & 200.5 \\ \hline
50 & 84.2 & 207.1 & 200.3 \\ \hline
60 & 89.4 & 207.0 & 200.3 \\ \hline
70 & 94.5 & 207.0 & 201.2 \\ \hline
80 & 99.5 & 207.5 & 202.0 \\ \hline
90 & 104.2 & 208.0 & 203.0 \\ \hline
100 & 109.5 & 209.5 & 204.0 \\ \hline
110 & 114.0 & 211.0 & 206.0 \\ \hline
120 & 120.0 & 213.0 & 208.2 \\ \hline
130 & 124.1 & 215.5 & 210.5 \\ \hline
140 & 129.2 & 218.5 & 213.5 \\ \hline
\end{longtable}

Программная обработка данных позволила рассчитать главные показатели преломления:
\begin{itemize}
    \item Для обыкновенной волны: \(n_o = 1.648 \pm 0.003\), \((\varepsilon = 0.18\%)\).
    \item Для необыкновенной волны: \(n_e = 1.477 \pm 0.005\), \((\varepsilon = 0.27\%)\).
\end{itemize}

На основе измеренных углов рассчитаны главные показатели преломления и построены графики зависимостей.  

\begin{figure}[H]
\centering
\begin{minipage}[b]{0.48\textwidth}
\centering
\includegraphics[width=\textwidth]{ordinary_refraction}
\caption{Зависимость \(n_o\) от угла падения \(\theta\)}
\label{fig:ordinary}
\end{minipage}
\hfill
\begin{minipage}[b]{0.48\textwidth}
\centering
\includegraphics[width=\textwidth]{extraordinary_refraction}
\caption{Зависимость \(n_e^2\) от \(\sin^2\theta\)}
\label{fig:extraordinary}
\end{minipage}
\end{figure}

Наблюдаемая слабая зависимость показателя преломления обыкновенной волны (\(n_o\)) от угла падения \(\theta\) (рис. \ref{fig:ordinary}) несколько не соответствует теоретическим ожиданиям для изотропных сред. Данный линейный тренд может быть связан с погрешностью юстировки: смещение призмы или лазера относительно оптической оси приводит к кажущемуся изменению \(n_o\) с ростом \(\theta\). Несмотря на это, усреднённое значение \(n_o = 1.648\) остаётся близким к теоретическому (\(n_{o,\text{теор}} = 1.65\)), что подтверждает корректность методики измерений.  

Для необыкновенной волны (\(n_e\)) экспериментальные данные демонстрируют чёткую линейную зависимость \(n_e^2\) от \(\sin^2\theta\) (рис. \ref{fig:extraordinary}), что согласуется с моделью анизотропного кристалла. Такое поведение объясняется зависимостью показателя преломления от ориентации вектора поляризации относительно оптической оси призмы. Полученные результаты подтверждают, что материал призмы обладает ярко выраженной анизотропией, характерной для одноосных кристаллов.  

\subsubsection*{Статистические параметры}  
\begin{itemize}
    \item Величина \(\chi^2\) для модели обыкновенной волны: $\chi^2_o = 34.93$.
    \item Величина \(\chi^2\) для модели необыкновенной волны: $\chi^2_e = 5.343$.
\end{itemize}

\subsection{Определение углов полного внутреннего отражения}  
Измерения проведены для обыкновенной и необыкновенной волн при нормальном падении света на призму. Результаты занесены в таблицу \ref{tab:critical_reflection}.

\begin{longtable}[c]{|c|c|c|}
\caption{Углы полного внутреннего отражения}
\label{tab:critical_reflection}\\
\hline
\textbf{Тип волны} & \textbf{Угол падения, \(^\circ\)} & \textbf{Отсчёт по риске, \(^\circ\)} \\ \hline 
\endfirsthead
%
\multicolumn{3}{c}{{\tablename} \thetable\ — продолжение} \\
\hline
\textbf{Тип волны} & \textbf{Угол падения, \(^\circ\)} & \textbf{Отсчёт по риске, \(^\circ\)} \\ \hline 
\endhead
%
Обыкновенная & \(230 \pm 3\) & \(59 \pm 0.5\) \\ \hline
Необыкновенная & \(225 \pm 4\) & \(51 \pm 0.5\) \\ \hline
\end{longtable}

\subsubsection*{Объяснение результатов}  
\begin{itemize}
    \item \textbf{Обыкновенная волна} (вертикальная поляризация) испытывает полное внутреннее отражение при угле падения \(230^\circ\), что соответствует отсчёту по риске \(59^\circ\). Это связано с превышением критического угла для данной поляризации.
    \item \textbf{Необыкновенная волна} проходит через призму, преломляясь вдоль гипотенузы, при угле \(225^\circ\) (отсчёт по риске \(51^\circ\)), так как её показатель преломления меньше.
\end{itemize}

\subsubsection*{Применение призмы Глана}  
В призме Глана два прямоугольных кристалла разделены воздушным зазором. При угле падения, превышающем критический для обыкновенной волны:
\begin{itemize}
    \item Обыкновенная волна отражается от границы раздела и поглощается.
    \item Необыкновенная волна проходит через призму, сохраняя горизонтальную поляризацию.
\end{itemize}

\begin{figure}[H]
\centering
\includegraphics[width=0.7\textwidth]{glan_prism_polarization}
\caption{Поляризация света в призме Глана}
\label{fig:glan_polarization}
\end{figure}

